% Generated draft: requires mathematical and source review.
\documentclass[11pt]{article}
\usepackage[margin=25mm]{geometry}
\usepackage{fontspec}
\setmainfont{Latin Modern Roman}
\usepackage{amsmath,amssymb,mathtools,unicode-math}
\setmathfont{Latin Modern Math}
\usepackage{enumitem,needspace,xurl,hyperref}
\setlength{\parindent}{0pt}
\setlength{\parskip}{5pt}
\setlength{\emergencystretch}{3em}
\newcommand{\N}{\mathbb{N}}
\newcommand{\Z}{\mathbb{Z}}
\newcommand{\Q}{\mathbb{Q}}
\newcommand{\R}{\mathbb{R}}
\newcommand{\C}{\mathbb{C}}
\newcommand{\F}{\mathbb{F}}
\newcommand{\A}{\mathbb{A}}
\newcommand{\PP}{\mathbb{P}}
\newcommand{\E}{\mathbb{E}}
\newcommand{\eps}{\varepsilon}
\newcommand{\dd}{\,\mathrm{d}}
\begin{document}
\section*{UnsolvedMath ID: 30004610}
\textbf{Source problem code:} catalogue entry
\subsection{Definitions and mathematical statement}
Fix a field \(K\). Let \(V\) be a finite-dimensional \(K\)-vector space. A central hyperplane arrangement \(\mathcal A\) in \(V\) is a finite set of codimension-one linear subspaces of \(V\); it is essential if
\[
\bigcap_{H\in\mathcal A}H=\{0\}.
\]
Here the standard finite-arrangement convention is used; the supplied record does not explicitly specify finiteness or discuss infinite arrangements.

The intersection lattice is
\[
L(\mathcal A)=\left\{\bigcap_{H\in\mathcal B}H:\mathcal B\subseteq\mathcal A\right\},
\]
where the empty intersection is \(V\), and the order is reverse inclusion: \(X\leq Y\) means \(X\supseteq Y\).

Put \(S_V=\operatorname{Sym}_K(V^*)\). For each \(H\in\mathcal A\), choose a nonzero linear form \(\alpha_H\in V^*\) with \(H=\ker\alpha_H\). Define the module of logarithmic derivations by
\[
D(\mathcal A)=\left\{\theta\in\operatorname{Der}_K(S_V):\theta(\alpha_H)\in\alpha_H S_V\text{ for every }H\in\mathcal A\right\}.
\]
This definition is independent of the choices of defining forms. The arrangement is free if \(D(\mathcal A)\) is a free \(S_V\)-module.

\textbf{Question (Terao's conjecture, as supplied).} For this fixed field \(K\), is it true that for every pair of finite central essential hyperplane arrangements \(\mathcal A\subseteq V\) and \(\mathcal B\subseteq W\), in finite-dimensional \(K\)-vector spaces, an isomorphism of intersection lattices implies
\[
L(\mathcal A)\cong L(\mathcal B)
\quad\Longrightarrow\quad
\bigl(\mathcal A\text{ is free}\ \Longleftrightarrow\ \mathcal B\text{ is free}\bigr)?
\]
No restriction on the characteristic of \(K\) is supplied. Thus \(K\) remains a fixed parameter; the supplied literature-status assessment is not taken as verification of the status for every possible field.
\subsection{Short English statement}
Terao's conjecture asks whether, over the same fixed field \(K\), two central essential hyperplane arrangements with isomorphic intersection lattices must either both be free or both fail to be free. The supplied record labels the question open, including in dimension three; this is a supplied dated assessment requiring independent verification.
\subsection{Sources}
\begin{itemize}
\item Local UnsolvedMath registry: supplied complete record with permanent ID \(30004610\), source problem code \texttt{catalogue entry}, in the Oberwolfach Reports Open Problems collection. Project registry locations are \texttt{data/unsolvedmath/problems.json} and the supplementary \texttt{catalog/problems.json}; their contents have not been independently inspected for this draft. The record names Terao's conjecture, gives proposed year 2021, and supplies no separate proposer.
\item Supplied source report: \emph{Logarithmic Vector Fields and Freeness of Divisors and Arrangements: New perspectives and applications} (2022), Oberwolfach Reports, pp.\ 225--300, DOI \texttt{10.4171/owr/2021/5}. \url{https://doi.org/10.4171/owr/2021/5}. The literature-source list labels the same reference \emph{Oberwolfach Report 5/2021}. The background describes these reports as documenting mathematical workshops and their open-problem discussions.
\item Supplied literature reference: Pagaria, \emph{Journal of the London Mathematical Society}. \url{https://doi.org/10.1112/jlms.12829}. No article title or more detailed bibliographic information is supplied.
\item Literature assessment dated exactly \texttt{2026-08-22}: the supplied curation labels Terao's conjecture open, including in dimension three, and lists the preceding Pagaria and Oberwolfach references as checked. This is supplied dated curation derived from the OWR record, not independently verified literature review; time-sensitive status and cited claims require independent verification.
\end{itemize}
\end{document}
